\cotefile{157-5}
\scanpage{558}
\mspage{469}
\unsure{dit|de} $N=\{\varnothing\}$, on voit que la condition supplémentaire
\deleted{\unsure{de sa|de la}} stabilité sur $\unsure{\underline T|\overline{\underline I}|\overline{\underline J}}$ n'est pas
superflue. (\uline{NB} Elle sert à assurer que
$\underline S$ est $L_\pi$-stable dans $M$ pour $\pi$ donné.)

4.16.3. Questions \unsure{apparentées} : j'avais espéré
que pour un foncteur accessible entre
catégories accessibles
\[
f:\unsure{N}\mathrel{\unsure{-|\longrightarrow}}M,
\]
la $\unsure{\leftrightarrows|\rightleftarrows}$-image essentielle, formée des $y\in\operatorname{Ob}M$
qui sont isomorphes à un facteur direct
d'un objet de la forme $f(x)$, où $x\in\operatorname{Ob}N$,
serait accessible dans $N$. L'exemple
ci-dessus (nullement \unsure{éliminatoire} !) de
la sous-catégorie $M_z$ de $M$ (formée des
$x\in\operatorname{Ob}M$ tels que $\operatorname{Hom}_M(x,z)\ne\varnothing$) montre
qu'il n'en est rien. En effet, $M_z$ est
l'image essentielle, et aussi l'image $\leftrightarrows$-essentielle,
du foncteur
\[
M_{/z}\longrightarrow M,
\]
où $M_{/z}$ est accessible et le foncteur précédent
accessible, si $M$ est \unsure{accessible} (\unsure{vérification}
par \unsure{AQT}, cf plus \deleted{\unsure{bas}} énoncé plus général ?).
% archive sheet number: 557
\scanpage{559}
\mspage{470}
On a l'impression que la raison \unsure{pour} laquelle
on a de tels contre-exemples, c'est que
l'on a des $\unsure{\leftrightarrows|\rightleftarrows}$-images essentielles qui ne sont
\deleted{\unsure{satisfont à}} $L_\pi$-stables dans $M$ pour aucun
cardinal $\pi\in\unsure{\mathcal U|U}$. Cela amène à \unsure{modifier}
la conjecture ainsi : soit $\pi$ un cardinal
tel que $M$ \unsure{satisfasse|satisfait} à $L_\pi$, et soit
$\underline S(\pi)$ la plus petite sous-catégorie \unsure{str. pl|str. ple} de
$M$, $L_\pi$-stable dans $M$, et qui contienne
l'image ess. de $f$. (On pourrait appeler
$\underline S(\pi)$ la $L_\pi$-image essentielle.) Alors $\underline S(\pi)$
est accessible dans $M$. De plus, la famille
(décroissante) des $\underline S(\pi)$, pour $\pi$ variable,
est stationnaire, i.e il existe un cardinal
$\pi$ tel que $\pi'\geq\pi$ implique $\underline S(\pi')=\underline S(\pi)$.
Il est clair \inserted{alors} que $\underline S(\pi)$ est \unsure{la|le} plus \unsure{petite|petit}
sous-catégorie de $M$ qui soit accessible dans $M$,
et qui contienne l'image essentielle de
$f$. La même question se pose pour
un foncteur $g^{\circ}:N^{\circ}\longrightarrow M$, où $g:N\mathrel{\unsure{-|\longrightarrow}}M^{\circ}$
est un foncteur accessible, $N$ et $M$ accessibles.
Je conjecture que la réponse est encore
\unsure{affirmative}.
% archive sheet number: 558
\scanpage{560}
\mspage{471}
Soit $M$ une catégorie satisfaisant la condition
$L$, et admettant une petite sous-catégorie
pleine strictement génératrice $C$, formée
d'objets accessibles. On a donc un
foncteur pl. fid. et accessible
\[
M\hookrightarrow C^{\wedge},
\]
$M$ \inserted{$L_\pi$-}stable dans $C^{\wedge}$ pour \deleted{\unsure{tout}} $\pi$ convenable.
Si $\overline M$ est l'image essentielle, de sorte
que $M\approx\overline M$, c'est donc une ss-catégorie
str. pleine $L_\pi$-stable dans $C^{\wedge}$. Donc on
s'attend à ce qu'elle soit accessible
dans $C^{\wedge}$ (si la situation décrite est
``logiquement accessible''), \deleted{\unsure{donc enfin}}
\deleted{\unsure{et}} \unsure{revient} au $=$ que $M$ soit accessible.
(\unsure{Comparer} 4.8.1 p. 187). Mais il ne doit
pas être difficile$^{*}$, en procédant à un
\unsure{nb|nbre} $\pi_{\unsure{\mathcal U|U}}$ de choix (où $\pi_{\mathcal U}=\operatorname{card}\mathcal U$), de
\unsure{fabriquer|trouver} quand même des exemples (``logiquement
\unsure{inaccessibles|accessibles}'') d'une sous-catégorie
str. pleine $M$ de $C^{\wedge}$ qui soit $L_\pi$-stable,
mais non accessible dans $C^{\wedge}$. Et $M$
% slanted
\margin{OK en prenant un \uline{\unsure{fermé}} \unsure{(image} accessible dans $M$ d'une catégorie accessible $M$, \unsure{qui se plonge} dans un $C^{\wedge}$. Mais \unsure{cet} $\underline S$ n'a \uline{pas} une \unsure{petite catégorie|catégorie petite} str. génératrice !}
% slanted
\margin{Ce n'est pas si clair, il faudrait pour cela que le foncteur d'injection ait un adjoint à \unsure{gauche|gche} !}
admet donc une petite sous-catégorie str.
génératrice \inserted{(formée d'objets accessibles$^{*}$)} et satisfait à $L_\pi$, sans pourtant
être accessible.
% archive sheet number: 559
\scanpage{561}
\mspage{472}
% slanted
\margin{\unsure{descendre} après 4.16.5.}
4.16.4. Soit $M$ une catégorie telle que $M$
et $M^{\circ}$ soient accessibles. Je conjecture
que alors $M$ est ess. petite. Plus
\unsure{généralement}, supposons $M$ accessible, et
que $M$ admette une petite sous-catégorie
qui soit \uline{strictement} \uuline{co}génératrice. Alors
$M$ est-elle ess. petite ? Plus généralement,
si $M$ admet une petite sous-catégorie
str. génératrice, et aussi une petite
sous-catégorie str. cogénératrice, i.e si
$M$ et $M^{\circ}$ sont toutes deux équivalentes
à des sous-catégories pleines d'une
$C^{\wedge}$, avec $C$ petite, alors $M$ est-elle
ess. petite ?

4.16.5. Soit $I$ une petite catégorie,
$I_{\unsure{\cdot|\bullet}}$ le cône droit sur $I$ (obtenu en rajoutant
un objet final $e$), $M$ une catégorie accessible,
\[
\underline{\operatorname{Hom}}_{\square}(I_{\cdot},M)\subset\underline{\operatorname{Hom}}(I_{\cdot},M)
\]
la sous-catégorie str. pleine de \deleted{$\operatorname{Hom}$} la catégorie
(accessible) $\underline{\operatorname{Hom}}(I_{\cdot},M)$, formée des foncteurs
$f$ qui font de $f(e)$ la limite inductive
dans $M$ des $f(i)$ ($i\in\operatorname{Ob}I$). \deleted{Je conjecture} \inserted{Cette ss-catégorie}
est $L_\pi$-stable pour $\pi$ grand, \unsure{donc} il y a lieu
% archive sheet number: 560
\scanpage{562}
\mspage{473}
de conjecturer qu'elle est accessible dans
$\underline{\operatorname{Hom}}(I_{\unsure{\cdot|\bullet}},M)$. [Je rappelle qu'on a vu
ds. 4.14. que c'est OK si les $\underline{\lim}$ de type
$I$ sont représentables dans $M$. Et qu'une
réponse affirmative, dans le cas où
\diagram{\begin{tikzcd}[ampersand replacement=\&,column sep=small,row sep=small]
\& \& 1\\
{I={}}\& 0\arrow[ur]\arrow[dr] \& {}\\
\& \& 2
\end{tikzcd}
\quad donc\quad
\begin{tikzcd}[ampersand replacement=\&,column sep=small,row sep=small]
\& \& 1\arrow[dr] \& {}\\
{I_{\cdot}={}}\& 0\arrow[ur]\arrow[dr] \& {} \& 3\\
\& \& 2\arrow[ur] \& {}
\end{tikzcd}}
, implique
que la sous-catégorie $\underline{\unsure{\operatorname{Ep}|\operatorname{Ef}}}(M)$ de $\underline{\operatorname{Fl}}(M)$
est accessible dans $\underline{\operatorname{Fl}}(M)$.]
\margin{\uline{NB} Si $M$ admet une petite sous-catégorie \uuline{co}génératrice, alors la question est justiciable de la conjecture 4.16.2, relative à un foncteur accessible $M\longrightarrow N^{\circ}$, ici $\underline{\operatorname{Hom}}(I_{\cdot},M)\longrightarrow\underline{\operatorname{Hom}}(I^{\circ}\times C,\operatorname{Ens})^{\circ}$.}

4.16.6. Il reste la question de ``l'accessibilité''
relative d'une catégorie $\unsure{\mathcal E|E}$ sur une
petite catégorie $C$. Je pense à présent,
grâce à 4.15.9, et son cor. 4.15.10, que l'énoncé suivant doit ``tomber'' par \unsure{AQT} :

\uline{Proposition} 4.16.6.1. Soit $f:\mathcal E\mathrel{\unsure{-|\longrightarrow}}C$ un
foncteur, avec $C$ petite catégorie \unsure{karoubienne}.
On suppose que $f$ satisfait la condition
$H'=H(1+2\unsure{+|\cdot}3)$ de 4.15.1 (page 415), \inserted{où il suffit $H2,H3$,} pour tte
flèche $\lambda$ \uline{non identique} de $C$. On suppose de
plus que pour \deleted{\unsure{toute flèche sous-}}\inserted{une} catégorie $C'$
\unsure{$\subseteq$|$C'$} de $C$ d'un des \inserted{\unsure{deux}} types qu'on va décrire, et un \inserted{foncteur $C'\longrightarrow C$,}
$\mathcal E'=\mathcal E\times_C C'$ soit accessible et $\mathcal E'\longrightarrow C'$ accessible\unsure{.|:}
% archive sheet number: 561
\scanpage{563}
\mspage{474}
\deleted{1°) Catégorie définie par un objet $x$ de $C$, un
projecteur $p$ de $x$ tel que $p\ne\operatorname{id}_x$, un
objet $x^{\unsure{\leftrightarrows|\natural}}$ et un couple \deleted{\unsure{de}} de flèches
\diagram{\begin{tikzcd}[ampersand replacement=\&]
x^{\leftrightarrows}\arrow[r,shift left,"\sigma"]\&x\arrow[l,shift left,"\rho"]
\end{tikzcd}}
faisant de $x^{\leftrightarrows}$ une image de $p$,
enfin une flèche $x\xrightarrow{u}y$ telle que
$up=u$, et la flèche composée}

[On prend pour $C'$ la catégorie
\diagram{\small\(\displaystyle \unsure{Q|\mathcal Q}=\left(\vcenter{\hbox{\begin{tikzcd}[ampersand replacement=\&,column sep=large,row sep=large]
{\unsure{\xi|\zeta}_0=\xi^{\leftrightarrows}}\arrow[r,shift left,"\sigma"]\arrow[dr,"u_0"']\&\xi\arrow[l,shift left,"\rho"]\arrow[loop right,"p"]\arrow[d,"u=u_0\rho"]\\
{}\& y
\end{tikzcd}}}\quad
\begin{array}{ll}
\rho\sigma=\operatorname{id}_{\deleted{\unsure{\xi|x}}\unsure{\xi_0}}&\\
\sigma\rho=p&\text{d'où }\left|\begin{aligned}p^2&=p\\p\sigma&=\sigma\\\rho p&=\rho\end{aligned}\right.\\
u=u_0\rho&\text{d'où }\left|\begin{aligned}up&=u\\u\sigma&=u_0\end{aligned}\right.\\
\deleted{\unsure{u\rho=}}&
\end{array}\right)\)}
(trois objets $\xi,\xi_0,y$ et cinq flèches non identiques
$\sigma,\rho,\unsure{p},u_0,u$, liées par les relations ci-dessus),
ou la catégorie (isomorphe à $Q^{\circ}$)
\diagram{\small\(\displaystyle Q'=\left(\vcenter{\hbox{\begin{tikzcd}[ampersand replacement=\&,column sep=large,row sep=large]
\xi_0\arrow[r,shift left,"\sigma"]\&\xi\arrow[l,shift left,"\rho"]\arrow[loop right,"p"]\\
{}\& y\arrow[u,"u"']\arrow[ul,"u_0"]
\end{tikzcd}}}\quad
\begin{array}{l}
\left.\begin{aligned}\rho\sigma&=\operatorname{id}_{\xi_0}\\\sigma\rho&=p\end{aligned}\right\}
\Rightarrow\left|\begin{aligned}p^2&=p\\p\sigma&=\sigma\\\rho p&=\rho\end{aligned}\right.\\
u=\sigma u_0\ \Rightarrow\left|\begin{aligned}pu&=u\\\rho u&=u_0\end{aligned}\right.
\end{array}\right)\)}
On \deleted{prend des}\inserted{se limite aux} foncteurs $Q\longrightarrow C$, $Q'\mathrel{\unsure{-|\longrightarrow}}C$
tels que$^{\unsure{*|!}}$ $p$ soit transformé en une
flèche \uline{non identique} de $C^{*}$ (donc la
% slanted
\margin{je ne suis pas sûr si c'est raisonnable --- à vérifier avec \unsure{soin} !}
condition relative à $Q,Q'$ devient vide, si
$C$ est sans projecteurs).]

Sous ces conditions, $\unsure{\mathcal E|E}$ est accessible, et
$f:\mathcal E\mathrel{-}C$ est accessible.
% archive sheet number: 562
\scanpage{564}
\mspage{475}
On trouve \unsure{alors|donc} \deleted{\unsure{comme}} le

\uline{Corollaire} 4.16.6.2. Soit $f:\unsure{\mathcal E|E}\mathrel{\unsure{-|\longrightarrow}}C$ un
foncteur, conditions équivalentes

a) Pour tte \deleted{\unsure{petite}} catégorie $C'$ \uline{petite} et
\uline{karoubienne} et tout foncteur $C'\mathrel{-}C$,
$\mathcal E'=\mathcal E\times_C C'$ est accessible, et le foncteur
$\mathcal E'\longrightarrow C'$ est accessible.

b) Comme a), mais $C'$ étant une
des \deleted{\unsure{cat.}} trois catégories \deleted{$\unsure{C|Q}$} $\Delta^1,\unsure{Q|\mathcal Q},Q'$
(où $Q,Q'$ sont définies dans 4.15.6.1 ci-dessus).

b$'$) (Si $C$ est sans projecteurs) Comme b),
mais en se limitant au cas $C'=\Delta^1$
(à l'exclusion de $Q,Q'$)
% slanted
\margin{? à vérifier avec soin !}

\uline{Corollaire} \unsure{4.15.6.3.|4.16.6.3.} Supposons $\deleted{\unsure{g|f}}f:\mathcal E\mathrel{-}C$
fibrant ou cofibrant. Alors les conditions
équivalentes précédentes impliquent que

1°) les $\mathcal E_s$ sont accessibles ($s\in\operatorname{Ob}C$) et

2°) pour tte flèche $\lambda:s\mathrel{-}t$ de $C$,
le foncteur $\lambda^{\unsure{*|\alpha}}:\mathcal E_t\mathrel{-}\mathcal E_s$ (resp $\lambda_{\unsure{*|<}}:\mathcal E_s\mathrel{-}\mathcal E_t$)
est accessible. \deleted{\unsure{Deuxièmement}} D'autre part,
l'une de ces deux conditions 1°, 2°) équivaut
à la validité de b) ci-dessus, dans le
cas de $C'=\Delta^{\unsure{1|'}}$, ou encore à la \unsure{solidité|validité}
% archive sheet number: 563
\scanpage{565}
\mspage{476}
de a) ci-dessus, quand $C'$ est \uline{sans projecteurs}.

\uline{Question} 4.16.6.4. Les conditions 1°) et 2°) du
\unsure{précédentes|précédent} impliquent-elles celles de 4.16.6.2 ?
En d'autres termes, si la catégorie base
$C$ est égale à $\unsure{Q|\mathcal Q}$ ou à $Q'$, est-ce que ces
conditions impliquent-elles que $\unsure{\mathcal E|E}$ est
accessible, et $f:\mathcal E\longrightarrow\deleted{\unsure{\mathcal E}}C$ accessible ?

\unsure{Ça|Sa} semble \unsure{idiot} de ne pas connaître
la réponse !

\S 4.16.6.5. J'ai envie d'appeler
catégorie \uline{accessible sur} $\deleted{\unsure{\mathcal E}}C$, une catégorie
$\mathcal E$ sur $C$ satisfaisant les conditions
équivalentes de \S 4.16.6.1 --- du moins
dans le cas où on suppose $C$ petite.
Pour que cette terminologie soit vraiment
\unsure{satisfaisante|satisfaisant}, il faudrait que l'on
sache prouver que ces conditions impliquent
ceci : pour toute catégorie accessible
$M$ et \unsure{tout} foncteur accessible $M\longrightarrow C$,
la catégorie $\deleted{\unsure{M|C}}\mathcal E\times_C M$ est accessible, et
\unsure{les|le} \unsure{foncteurs|foncteur} $\mathcal E\times_C M\mathrel{\unsure{-|\longrightarrow}}M$ est accessible.
% archive sheet number: 564
\scanpage{566}
\mspage{47\deleted{\unsure{7}}6 bis}
L'idée, c'est qu'on appellerait (sans hypothèse
de petitesse) \uline{catégorie accessible} $\unsure{\mathcal E|E}$ sur une
autre $C$, une catégorie ayant la propriété
qu'on vient de dire --- mais en supposant
\inserted{\unsure{sans doute|toutefois}} de plus que $C$ satisfasse à $L_\pi$, j'imagine :
pour tt foncteur accessible
$M\mathrel{\unsure{-|\longrightarrow}}C$, avec $M$ accessible,
$\mathcal E\times_C M$ est accessible et le foncteur
canonique $\mathcal E\times_C M\longrightarrow\mathcal E$ est accessible.

\deleted{4.17. \uline{Retour sur les \deleted{\unsure{opérations}}\inserted{\unsure{paratopos}} et les
opérations} $M\boxtimes N,\underline{\operatorname{Hom}}_{!}(M,N)$.}

\deleted{4.16.7. Retour sur 4.16.1 :

\uline{Théorème} 4.16.7.1. Soit $M$ une catégorie accessible.
Alors a) Toute sous-catégorie \uline{fermée} $\underline S$ de $M$
est accessible dans $M$, et b) L'ensemble
de ces sous-catégories est \uline{petit}.}

\deleted{\uline{Corollaire} 4.16.7.2. Soit $M$ une catégorie accessible.
Pour que toute sous-catégorie ouverte $\unsure{\underline T|\overline{\underline I}|\overline{\underline J}}$ de
$M$ soit accessible, il f. et s. qu'il en
soit ainsi des \deleted{\unsure{catég}} sous-catégories $M_z$
considérées dans 4.16.1., i.e. (pour $z\in\operatorname{Ob}M$)}
% archive sheet number: 565
\scanpage{567}
\mspage{478}
\deleted{de la \inserted{sous-}catégorie \deleted{\unsure{str. pleine}} ouverte de $M$
formée des $x\in\operatorname{Ob}M$ tels que $\operatorname{Hom}_M(x,z)\ne\varnothing$.}

\deleted{En effet, toute \deleted{\unsure{ouverte est}} sous-catégorie ouverte
\inserted{$\unsure{\underline T|\overline{\underline I}|\overline{\underline J}}$} est une réunion des sous-catégories ouvertes
de la forme $M_z\unsure{.}$ (pour $z\in\operatorname{Ob}\underline T$). Comme
l'ensemble des \deleted{\unsure{ces}} sous-catégories ouvertes
de $M$ est petit (celui des sous-catégories
fermées l'étant), il s'ensuit que toute
sous-catégorie ouverte est une réunion
d'une famille \uline{petite} de sous-catégories
$M_z$. Si ces dernières sont accessibles
dans $M$, il en est donc de $=$ de leurs
réunions en vertu de 4.13\ldots, d'où le
corollaire.}

\deleted{Démonstration de 4.16.7.1. Je vais d'abord
prouver b). Notons le

\uline{Lemme} 4.16.7.3. Soit \deleted{\unsure{$C$}} $f:M$\inserted{$\longrightarrow N$} \unsure{une} \deleted{catégorie,}
\deleted{\unsure{$N$ une sous}} un foncteur pl. fidèle entre
deux catégories. Alors l'application
\[
\alpha:N'\longmapsto f^{-1}(M'):\unsure{\mathcal F|F}(N)\mathrel{\unsure{-|\longrightarrow}}\mathcal F(M)
\]
est surjective, où $\mathcal F(N)$ est \unsure{l'ens.} des
\deleted{\unsure{parties}}\inserted{ss-cat.} fermées \deleted{\unsure{du}}de $\deleted{\unsure{M}}N$. De façon plus
précise, considérons l'application \deleted{\unsure{$\beta:\mathcal F(M)\longrightarrow\mathcal F(N)$}, $\unsure{M'\mapsto N'}$, où $\unsure{N'=f(M')|N'=f_!M'}$ \inserted{$\unsure{f_!M'|f(M')}$} \unsure{ensemble} $\unsure{f(\operatorname{Ob}M)}$}}
% archive sheet number: 566
\scanpage{568}
\mspage{477}
\uline{4.1\deleted{\unsure{5}}6.7.} \uline{Retour sur les sous-catégories fermées d'une
catégorie accessible} (4.16.1).

\uline{Proposition} 4.16.7.1. \unsure{Soient|Soit} $M$ une catégorie
accessible, \deleted{$\pi_0$ un cardinal \inserted{\unsure{infini}} tel que $M$ satisfasse
à $L_{\pi_0}$,} $(\underline S_i)_{i\in I}$ une famille de sous-catégories
\inserted{fermées} accessibles dans $M$\deleted{, et $L_{\pi_0}$-stables
dans $M$}. On ne suppose pas que \unsure{card} $I$ soit
petit. Soit $\underline S$ la sous-catégorie \inserted{str.} pleine
de $M$ telle que $\deleted{\unsure{\underline S}}\operatorname{Ob}\underline S=\bigcup_{i\in I}\operatorname{Ob}\underline S_i$.
Alors pour que $\underline S$ soit accessible dans
$M$, il f. et s. qu'il existe une
\uline{petite} partie $J\subset I$, telle que
$\operatorname{Ob}\underline S=\bigcup_{i\in J}\operatorname{Ob}\underline S_i$.
\inserted{(En particulier, si $I$ est ordonné et $i\mapsto\underline S_i$ \deleted{\unsure{str.}} croissante, \inserted{et \unsure{gd|grd} devant \unsure{tt|ttes} petites parties,} il f. et s. qu'il existe un $i\in I$ tel que \deleted{\unsure{$\underline S_i$}} $\underline S=\underline S_i$, i.e tel que $j\geq i$ implique $\underline S_j=\underline S_i$.)}

C'est évidemment suffisant, en
vertu de 4.13.\ldots\ Prouvons \unsure{inversement|l'assertion} que
si $\underline S$ est accessible \unsure{dans}\inserted{$M$}, alors la condition
précédente est satisfaite. Soit $\pi$ un cardinal
$\geq\pi_0$ tel que $\underline S$ soit $\pi$-accessible dans
$M$, donc $\underline S\xleftarrow{\approx}\operatorname{Ind}_\pi\underline S(\pi)$, où $\underline S(\pi)=\underline S\cap M_\pi$.
% archive sheet number: 567
\scanpage{569}
\mspage{478}
\deleted{Pour tout} Soit $\underline S(\pi)_0$ une sous-catégorie réduite
de $\underline S(\pi)$, telle que $\underline S(\pi)_0\hookrightarrow\underline S(\pi)$ soit
une équivalence, de sorte que $\underline S(\pi)_0$ est
petite. Pour tout $x\in\operatorname{Ob}\underline S(\pi)_0$, soit
$i_x\in I$ \deleted{\unsure{un des indices qui}}\inserted{tel que} $x\in\operatorname{Ob}\underline S_{i_x}$. \unsure{Soit donc|Soit}
$J$ l'ens de tous les $i_x$, qui est
une petite partie de $I$, et soit
$\underline S_J$ la sous-catégorie \inserted{str. pl.} de $M$ ayant
pour \deleted{\unsure{ens. d'}} ens. d'objets $\bigcup_{j\in J}\operatorname{Ob}\underline S_j$.

\deleted{\unsure{On a}} Donc $\underline S(\pi)_0$ est contenue
dans $\underline S_J$, donc aussi $\underline S(\pi)$ puisque
$\underline S_{\deleted{\unsure{i}}J}$ est strictement pleine :
\[
\underline S(\pi)\subset\underline S_J.
\]
D'ailleurs $\underline S_J$ est accessible dans $M$
(loc. cit), et est fermée dans $M$, comme
réunion de sous-catégories fermées, \uline{a fortiori}
\uline{elle satisfait à} $L_\pi$. Donc, comme $\underline S(\pi)\subset\underline S_J$,
le foncteur \inserted{(\unsure{ess.} surjectif)} $\underline{\operatorname{Ind}}_\pi(\underline S(\pi))\longrightarrow\underline S$ se
factorise par $\underline S_J$. Donc $\underline S_J\hookrightarrow\underline S$ est
ess. surjectif, et comme $\underline S_J$ est str.
pleine, on en conclut $\underline S_J=\underline S$, \unsure{qed|cqfd}.
% archive sheet number: 568
\scanpage{570}
\mspage{479}
\uline{Proposition} 4.16.7.2. \unsure{Soient|Soit} $M$ une catégorie
accessible, $x$ un objet de $M$, et $ {}_xM$ la
sous-catégorie fermée de $M$ engendrée
par $x$, formée des $y\in\operatorname{Ob}M$ tels que
$\operatorname{Hom}(x,y)\ne\varnothing$, i.e. \deleted{\unsure{de}} l'image essentielle
de $ {}_{x\backslash}M\mathrel{\unsure{-|\longrightarrow}}M$. Alors $ {}_xM$ est une
sous-catégorie de $M$ accessible dans $M$.

En effet, c'est l'image inverse de la
sous-catégorie $\operatorname{Ens}^{*}$ accessible dans $\unsure{\operatorname{En}|\operatorname{Ens}}$,
par le foncteur représentable (donc
accessible)
\[
h_x:y\longmapsto\operatorname{Hom}(\unsure{xy|x,y}):M\mathrel{-}(\operatorname{Ens}),
\]
donc est accessible \inserted{dans $M$} en vertu de 4.13.\ldots

\uline{\deleted{Proposition}\inserted{Corollaire}} 4.16.7.3. Soit $M$ une catégorie
accessible, $\underline S$ une sous-catégorie fermée.
Pour que $\underline S$ soit accessible, il f. et il
s. qu'il existe une \uline{petite} partie
\inserted{$J$} de $\underline S$, telle que $\operatorname{Ob}\underline S=\bigcup_{x\in J}\deleted{\unsure{\mathcal E}}\operatorname{Ob}({}_xM)$.

En effet, on a
\[
\operatorname{Ob}\underline S=\bigcup_{x\in\operatorname{Ob}\underline S}\operatorname{Ob}({}_xM),
\]
i.e. $\underline S$ apparaît comme réunion des
% archive sheet number: 569
\scanpage{571}
\mspage{480}
de la famille (pas néc. petite) de sous-catégories
fermées $ {}_xM$ de $M$ (\unsure{indexée|indexées} par
$x\in\operatorname{Ob}\underline S$), \inserted{accessibles dans $M$ en vertu de 4.16.7.2,} donc notre assertion est un
cas particulier de 4.16.7.1.

On peut aussi, pour toute partie petite
$J$ de $\operatorname{Ob}\underline S$, considérer la sous-catégorie
fermée $\underline S_J$\inserted{$\subset\underline S$} de $\deleted{\unsure{\underline S}}M$ engendrée
par $J$. L'ens \inserted{$I$} des parties $J$ en
question, ordonné par inclusion, est
filtrant devant \deleted{\unsure{tout}} cardinal $\pi\in\unsure{\mathcal U|U}$, et
$\underline S$ est réunion croissante des $\underline S_J$.
D'autre part, la proposition nous dit
que $\underline S$ est accessible dans $M$ ssi
\deleted{\unsure{$J$}} il existe un $J\in I$, tel que $\underline S_J=\underline S$,
i.e. tel que $\underline S_J=\underline S_{J'}$ pour $J\unsure{\subseteq|\subset}J'$.

\uline{Proposition} 4.16.7.4. Soit $M$ une catégorie
accessible. Les conditions suivantes
sont équivalentes

a) Il existe une sous-catégorie $\underline S$ de $M$
fermée dans $M$, et non accessible dans $M$.

b) \deleted{Il existe \unsure{soit|des} \unsure{$\alpha$|objets} \unsure{premier|paramétrés} ordinal} \inserted{Soit $I$ l'ens des ordinaux $\alpha$}
tels que $\operatorname{card}\alpha$ \deleted{\unsure{$\leq$}} $\in\mathcal U$ \deleted{\unsure{; (donc}}\inserted{\unsure{i.e.}} des
% archive sheet number: 570
\scanpage{572}
\mspage{481}
ordinaux $\alpha$ \deleted{\unsure{$\in$}} qui sont $\in\unsure{\mathcal U|U}$. Donc $I=I_{\alpha_0}$, où
$\alpha_0$ est le plus petit ordinal qui n'appartient
pas à $\mathcal U$, et $I$ est grand devant tt
\inserted{petit} cardinal \deleted{\unsure{de}} \unsure{i.e.|de} devant tt cardinal $\in\mathcal U$.
% slanted
\margin{introduire $I=I_{\mathcal U}$ avant l'énoncé de la proposition !}
La condition envisagée est la suivante :

Il existe une suite transfinie
\[
\underline x=(x_\alpha)_{\alpha\in I}
\]
d'éléments de $\operatorname{Ob}M$, telle que l'on ait
\[
\operatorname{Hom}(x_{\deleted{\unsure{\alpha}}\beta},x_{\deleted{\unsure{\beta}}\alpha})=\varnothing\qquad\text{si }\beta<\deleted{\unsure{\beta}}\alpha.
\]

De façon précise, pour une telle suite
transfinie, soit $\underline S\deleted{\unsure{\subset}}\unsure{=}\,{}_{\underline x}M$ la sous-catégorie
fermée de \deleted{\unsure{$\underline S$}} $M$ engendrée par la
famille des $x_\alpha$, \inserted{i.e telle que $\operatorname{Ob}\underline S=\bigcup_{\alpha\in I}\operatorname{Ob}{}_{x_\alpha}M=\bigcup_{\alpha\in I}x_\alpha\backslash\operatorname{Ob}M$.}
Alors $ {}_{\underline x}M$ est une
sous-catégorie fermée non accessible dans $M$,
et toute \deleted{\unsure{telle}} sous-catégorie fermée non
accessible dans $M$ s'obtient de cette
façon.

\uline{Démonstration} Considérons une suite
$\underline x=(x_\alpha)_{\alpha\in I}$ comme dessus, \deleted{\unsure{et la}}
\deleted{En vertu de 4.16.7.1,} si $J$ est une
petite partie de $I$, \deleted{et si \unsure{$\alpha$}} il existe
un $\alpha\in I$ qui majore les $\beta\in J$,
donc $\operatorname{Hom}(x_\beta,x_\alpha)=\varnothing$ pour $\beta\in J$,
% archive sheet number: 571
